\documentclass[11pt]{amsart}
\usepackage[T1]{fontenc}
\usepackage{lmodern}
\usepackage[margin=1in]{geometry}
\usepackage{amsmath,amssymb,mathtools,booktabs}
\usepackage{microtype}
\usepackage[hidelinks]{hyperref}
\numberwithin{equation}{section}
\newcommand{\R}{\mathbb R}
\newcommand{\C}{\mathbb C}
\newcommand{\Tr}{\operatorname{Tr}}
\newcommand{\rank}{\operatorname{rank}}
\newcommand{\supp}{\operatorname{supp}}
\newcommand{\HS}{\mathrm{HS}}
\newcommand{\op}{\mathrm{op}}
\newcommand{\dd}{\,\mathrm d}
\newcommand{\Nd}{N^{\rm d}}
\newcommand{\ind}{\mathbf 1}
\newcommand{\osc}{\operatorname{osc}}
\title[Compensated Gram corrections and zero counts]{Compensated Gram corrections for zero counts of the Riemann zeta function}
\author{R. Arun Chandru}
\address{Department of Aerospace Engineering \\ Indian Institute of Science}
\email{arun.chandru@gmail.com}
\date{10 September 2026}
\subjclass[2020]{Primary 11M26; Secondary 15A42, 65G20}
\keywords{Riemann zeta function, Gram matrix, zero multiplicity, pair correlation}
\begin{document}
\begin{abstract}
We develop a compensated Gram-matrix correction for counting zeros of the
Riemann zeta function. A single averaged matrix witness combines overlapping
local windows without an outer-block loss. Its admissible local floor depends
on the spectral clipping threshold; this dependence permits a stronger
correction for the union of simple and critical-line zeros and for
low-multiplicity counts than for the simple-critical intersection.
Combining the finite argument with unconditional pair correlation and two
computer-assisted seven-point inequalities gives lower limiting proportions
exceeding $0.888554068993217696$ for simple-or-critical zero mass,
$0.837368508521536857$ for mass of multiplicity at most two, and
$0.939197203148497742$ for locations of multiplicity at most two among
distinct locations. Each improves an explicitly identified same-statistic
comparison. We also obtain a simple-critical bound, a total-distinct bound,
fixed higher-multiplicity-cutoff families, and qualified short-height
extensions. The analytic deductions and the finite certification rules are
given explicitly. Supplement information treats the numerical certificates
and a local thirteen-point reduction; that reduction is not used as a
whole-domain inequality or as a further zero-count improvement.
\end{abstract}
\maketitle
\tableofcontents

\section{Introduction}

The distinction between a zero and a zero location is essential in questions
about multiplicity. So is the distinction between being simple \emph{and}
lying on the critical line and being simple \emph{or} lying there. These
questions share analytic information, but they spend that information in
different ways. A method that improves one count need not improve another.
This work tries to make that distinction effective through a
finite Gram-matrix argument.

The starting point is the unconditional pair-correlation theorem of Baluyot,
Goldston, Suriajaya and Turnage-Butterbaugh \cite{BGST}. The recent arguments
of Alp\"oge--Furman and Lamzouri \cite{AF,Lamzouri} connect this analytic
input to the simultaneous simplicity and criticality of more than two thirds
of the zeros. Local Gram corrections, including computer-assisted weighted
overlap inequalities, subsequently appear in the public work of Ainta and
Shi \cite{Ainta,Shi}. Waddle's public low-multiplicity draft \cite{Waddle}
provides a further useful comparison: its scalar multiplicity accounting
can be strengthened by retaining the entire simple-real Gram block.

Our contribution is a finite transfer mechanism and its application to these
different counts. The main point is not to add unrelated local estimates to
a spectral bound. Instead, all overlapping windows are represented inside
one admissible matrix witness. The quadratic cost of that witness pays for
their overlap. A row-sum estimate then controls the windows that require
clipping. Since that control grows as the square of the clipping threshold,
a local floor that is too large for the simple-critical argument can still
be used for the union and multiplicity arguments.

There are two independent numerical inputs. One is Ainta's seven-point
cosine-kernel inequality, which we credit as prior work. The other uses an
explicit two-mode perturbation of the same density. The perturbation raises
the analytic energy, so its local improvement is useful only after both
effects have been included in the same calculation. We give the exact
formulas, rather than combining the energy of one kernel with the local
certificate of another. The computer-assisted inputs are whole-domain
inequalities, not sampled minima. Their finite acceptance rules and
reproducibility requirements are supplied in the accompanying supplement.

All results in this article concern ordinary nontrivial zeros of $\zeta$.
Neither the Riemann hypothesis nor a narrow-box assumption on the zeros is
used in the global conclusions. The short-height statements additionally
invoke the stated theorem of Wang \cite{Wang}, currently a preprint; the
global conclusions do not depend on it.

\subsection{Counts and main conclusions}

Write a nontrivial zero location as $\rho=\beta+i\gamma$, and write
$m_\rho$ for its multiplicity. All sums in this subsection run over distinct
locations with $0<\gamma\le X$. Define
\begin{align}
 N(X)&=\sum_\rho m_\rho,&
 n(X)&=\#\{\rho:m_\rho=1,\ \beta=\tfrac12\},\label{eq:counts1}\\
 \Nd(X)&=\sum_\rho1,&
 U(X)&=\sum_{m_\rho=1\ \text{or}\ \beta=1/2}m_\rho,\label{eq:counts2}\\
 M_{\le r}(X)&=\sum_{m_\rho\le r}m_\rho,&
 D_{\le r}(X)&=\sum_{m_\rho\le r}1.\label{eq:counts3}
\end{align}
Thus $U$ is a mass count of a union, whereas $n$ counts the intersection.
The denominator for $D_{\le r}$ below is $\Nd$, not $N$.

Here are the exact constants used throughout. Put
\begin{gather}
 z=\frac1{\sqrt2},\qquad R_0=\frac12+z\cot z,
 \qquad u_1=\frac{13}{500},\quad u_2=-\frac7{500},\label{eq:R0}\\
 \Delta=\sum_{j=1}^2\left(\frac12-\frac1{4\pi^2j^2}\right)u_j^2,
 \qquad R_1=R_0+\Delta,\label{eq:Delta}\\
 a_0=\frac{19}{5000},\quad d_0=\frac1{500},\qquad
 a_1=\frac7{400},\quad d_1=\frac3{250},\quad T_1=R_1+d_1,\label{eq:rowconstants}\\
 \Gamma=4+2\sqrt{2(1-a_1)},\qquad
 Z=\frac{\Gamma^2}{8}=3-a_1+2\sqrt{2(1-a_1)},\qquad
 \beta=T_1-1-a_1,\qquad E=\frac{\beta}{Z}.\label{eq:E}
\end{gather}
In particular $0<E<1/2$. The conclusions proved below are
\begin{align}
 \liminf_{X\to\infty}\frac{n(X)}{N(X)}
 &\ge C:=\frac{2-R_0-d_0}{1-a_0}
 >0.673058325315610967,\label{eq:headlineintersection}\\
 \liminf_{X\to\infty}\frac{\Nd(X)}{N(X)}
 &\ge\frac{1+C}{2}>0.836529162657805483,\label{eq:headlinedistinct}\\
 \liminf_{X\to\infty}\frac{U(X)}{N(X)}
 &\ge1-2E>0.888554068993217696,\label{eq:headlineunion}\\
 \liminf_{X\to\infty}\frac{M_{\le2}(X)}{N(X)}
 &\ge\frac{3-T_1}{2-a_1}>0.837368508521536857,\label{eq:headlinemass}\\
 \liminf_{X\to\infty}\frac{D_{\le2}(X)}{\Nd(X)}
 &\ge1-\frac{T_1-1-a_1}{8-2T_1-a_1}
 >0.939197203148497742.\label{eq:headlinelocations}
\end{align}
The strict decimal inequalities compare fixed real constants. They do not
assert that every finite-height ratio already exceeds the printed decimal.
For each fixed integer $r\ge3$, the same argument also gives
\begin{align}
 \liminf_{X\to\infty}\frac{M_{\le r}(X)}{N(X)}
 &\ge1-\frac{r+1}{r-1}E,\label{eq:highermass}\\
 \liminf_{X\to\infty}\frac{D_{\le r}(X)}{\Nd(X)}
 &\ge1-\frac{E}{r-1-rE}.\label{eq:higherlocations}
\end{align}
The word ``fixed'' matters: no uniform asymptotic assertion with a
height-dependent cutoff is intended.

\subsection{Comparison with earlier bounds}
\label{subsec:literaturecomparison}

The union, multiplicity-two mass, and multiplicity-two distinct-location
conclusions improve identified unconditional bounds for the same counting
statistics. The union bound improves Lamzouri's
Theorem~1.1 \cite{Lamzouri}; the two multiplicity-two bounds improve
the displayed formulas in Waddle's public research draft \cite{Waddle}.
The respective improvements, as differences of proportions, exceed
\[
 0.000934060819863356,\qquad
 0.001118156681831034,\qquad
 0.000469272388685036.
\]
Multiplying these differences by $100$ gives percentage-point gains.
The comparisons use the exact expressions, not subtraction of rounded
headline percentages.

To state the comparisons compactly, put
\begin{equation}\label{eq:literatureconstants}
 C_0=2-R_0,\qquad C_1=\frac{3-R_0}{2},\qquad
 B=(3-2\sqrt2)(R_0-1),\qquad C_2=1-2B.
\end{equation}
Thus $C_0=0.672500703679411\ldots$ is the Montgomery--Taylor
simple-critical constant and $C_1=0.836250351839705\ldots$ is the
corresponding distinct-zero constant. Alp\"oge and Furman
\cite[Theorem A]{AF} established these bounds unconditionally by an
inertia and rank--trace argument for a finite compression of Weil's
Hermitian form. Lamzouri \cite{Lamzouri} recovered them through a
Hilbert-space argument and obtained the union constant
\[
 C_2=\frac{1+2\sqrt2+2C_0}{3+2\sqrt2}
       =0.887620008173354\ldots.
\]

\begin{table}[ht]
\centering
\small
\setlength{\tabcolsep}{4pt}
\renewcommand{\arraystretch}{1.6}
\begin{tabular}{@{}p{.10\textwidth}p{.43\textwidth}p{.40\textwidth}@{}}
\hline
Ratio & Present lower bound & Named comparison\\
\hline
$U/N$
& $1-2E$\par $>0.888554068993217696$
& Lamzouri \cite{Lamzouri}:\par
  $C_2=0.887620008173354\ldots$\\
$M_{\le2}/N$
& $\dfrac{3-T_1}{2-a_1}$\par $>0.837368508521536857$
& Waddle \cite{Waddle}:\par
  $C_1=0.836250351839705\ldots$\\
$D_{\le2}/N^{\rm d}$
& $1-\dfrac{\beta}{8-2T_1-a_1}$\par $>0.939197203148497742$
& Waddle \cite{Waddle}:\par
  $1-\dfrac{R_0-1}{2(4-R_0)}$\par
  $=0.938727930759812\ldots$\\
$n/N$
& $\dfrac{2-R_0-d_0}{1-a_0}$\par $>0.673058325315610967$
& Devine \cite{Devine}:\par
  $673399/10^6=0.673399$\par
  (stated unconditional theorem)\\
$N^{\rm d}/N$
& $\dfrac{3-R_0-d_0-a_0}{2(1-a_0)}$\par $>0.836529162657805483$
& Alp\"oge--Furman and Lamzouri\par
  \cite{AF,Lamzouri}:\par
  $C_1=0.836250351839705\ldots$\\
\hline
\end{tabular}
\caption{Five separate global lower limiting proportions. Each exact
expression in the middle column bounds the corresponding lower limit
from below; the strict inequalities compare these expressions with
shorter decimals. All external statements listed here are unconditional.
Waddle's source is an unrefereed public draft and Devine's is a public
preprint. The last row is a comparison with the stated foundational
bounds, not a claim of an overall distinct-zero record.}
\label{tab:literaturecomparison}
\end{table}

Here $n$ counts simple critical zeros, whereas $U$ counts with
multiplicity the union of simplicity and critical-line location.
The denominators $N$ and $N^{\rm d}$ are different: $M_{\le r}/N$
is a fraction of zero mass, and $D_{\le r}/N^{\rm d}$ is a fraction
of distinct complex zero locations. Locations with the same imaginary
part are not identified. These conventions also match the external
statements in Table~\ref{tab:literaturecomparison}. A dyadic statement
for additive counts gives the corresponding cumulative lower limit by
summing sufficiently high dyadic pieces of $(0,X]$; the bounded initial
remainder disappears after division by the same denominator.

For every fixed integer $r\ge3$, Waddle's displayed mass and
distinct-location bounds are respectively
\[
 1-\frac{(r+1)B}{r-1},\qquad
 1-\frac{B}{r-1-rB}.
\]
Our corresponding formulas replace $B$ by $E$, and the proved
inequality $0<E<B<1/2$ makes both improvements strict. The mass
gain is $(r+1)(B-E)/(r-1)$; the location comparison follows from
the strict increase of $x/(r-1-rx)$ on the relevant interval.
This conclusion uses the retained multiplicity budget. It is not a
formal consequence of a union bound: a multiset supported on real
locations of multiplicity $r+1$ has full union mass and no mass at
multiplicity at most $r$.

The fixed-cutoff statement is an improvement over a specified family.
It does not assert the best available bound for every cutoff.
Simoni\v{c}, Trudgian and Turnage-Butterbaugh \cite[Theorems 4 and 7]{STT}
give unconditional exponential estimates for the number of locations
of exact multiplicity $j$. Weighting by $j$ and summing the tail gives
mass bounds tending to one as the cutoff grows, whereas the displayed
mass family here tends to $1-E<1$. Consequently the exponential-tail
bound is stronger for sufficiently large fixed cutoffs.

The simple-critical conclusion refines Ainta's public bound
\cite{Ainta}, whose exact expression is
\[
 \frac{1345000C_0-2680}{1340003}
       =0.673008527927779\ldots.
\]
It remains below Shi's stated research-draft candidate
$0.6733169771424713\ldots$ \cite{Shi} and Devine's stated
$0.673399$ \cite{Devine}. The latter preprint's numerical proof has
not been independently reconstructed here. These are external
comparison targets, not premises of our proof, and we claim no overall
simple-critical or distinct-zero record. Zeta Lab's public report
\cite{ZetaLab} also records an eight-point value
$0.67305298298962888\ldots$, distinguishing an interval-certified
local inequality from its Lean implication conditional on that
inequality. The public union extension \cite{Union} supplies another
explicit treatment of the distinction between union and intersection.

Stronger numbers with additional hypotheses must be compared separately.
The union bounds $8/9$ in \cite[Remark 2]{BGSTI} and
$(2+C_0)/3=0.890833567893\ldots$ discussed in \cite[Remark 1.2]{Lamzouri}
use a narrow-box assumption on the horizontal locations of all
sufficiently high dyadic zeros; they are not unconditional union
comparators. Under RH, Gon\c{c}alves, de Laat and Leijenhorst
\cite[Theorem 1]{GdLL} prove
$\liminf M_{\le2}/N\ge0.9614$. Their cutoff-three value $0.9787$
also carries the theorem's additional qualification concerning Maple
summation of specified series. RH-dependent pair-correlation estimates
in \cite{CGdL} include $0.6792$ for simple zeros and $0.8477$ for
distinct zeros. Devine's claims beyond $0.6792$ likewise require
additional density or horizontal-profile assumptions \cite{Devine}.
None of these numbers is an unconditional obstruction to the bounds
proved here.


\subsection{Relation to earlier methods}
\label{subsec:literaturenovelty}

The analytic framework and several finite ingredients have established
antecedents. Montgomery's pair-correlation theorem \cite{Mon73} and
the Montgomery--Taylor window reported in \cite{Mon75} underlie the
constant $R_0$; the associated extremal problem is treated in
\cite{CCLM}. Aryan \cite{Aryan} obtained an unconditional form of
the relevant Fej\'er-kernel calculation, and Baluyot, Goldston,
Suriajaya and Turnage-Butterbaugh \cite{BGST} proved the published
full-height pair-correlation formula used here. Lamzouri's
finite-multiset formulation and removal of the rational pair weight
\cite{Lamzouri} provide the direct analytic interface. Wang
\cite{Wang} supplies the separate short-height input, with support
strictly smaller than the fixed interval exponent. We prove no new
pair-correlation theorem and do not insert a height-dependent test
into a fixed-test error estimate.

Ainta \cite{Ainta} introduced the clipped spectral stability
correction and quantitative overlapping seven-point windows used in
this line of refinements. The Fourier perturbations of the cosine
density and nonuniform pair weights with distance-capacity constraints
in the inspected Ojha and Shi drafts \cite{Ojha,Shi} are also prior
work. Shi's threshold-one trace--energy envelope, pressure rescaling,
and two-certificate supporting-plane deduction are directly relevant
antecedents. In particular, neither weighted overlap nor removal of
the elementary threshold-one cap is claimed as a new idea here.
The arbitrary-parameter scalar multiplicity framework and its
uncorrected low-multiplicity consequences are credited to the
two-trace development and Waddle's displayed deductions
\cite{AF,Waddle}.

The finite contribution here is the use of one common averaged witness
to transfer a whole-domain local Gram inequality directly to the
clipped spectral defect while preserving its pressure floor. The
uniform coefficient estimate depends on the clipping threshold.
For the seven-point weights used here, its admissible cap is
$40t^2/2401$. This dependence is material: the new local floor
$a_1=7/400$ is unavailable through this cap at $t=1$, but is
admissible at $t=2$ and at the union threshold
$t=1+\sqrt{2(1-a_1)}$. Thus the stronger local certificate is matched
to the count being proved. Retaining the resulting simple-Gram
correction in the separate per-location multiplicity costs then
produces the mass and distinct-location bounds from that same
certificate. The elementary variational identity for clipping is
part of the proof; the gain comes from its common-witness application
and the uniform capacity estimate.

The required two-mode profile, its normalization and energy cost,
and the whole-domain seven-point inequality are established explicitly.
The local interval engine is an attributed adaptation of Ainta's
MIT-licensed implementation, using Arb ball arithmetic \cite{Johansson}.
The separate cosine row for the simple-critical and total-distinct
conclusions is Ainta's completed input. Our reproduced executions
and exact scalar calculations certify the stated finite numerical
premises; they do not replace the finite operator argument or the
passage to zero limits. 

Finally, bounded-potential versions of the common-witness argument
give additional finite inequalities, but the general telescoping
subaction idea already occurs in the public literature, including
\cite{Devine}. Partial larger-window certificates are identified by
their domains. They supply no global zero-count conclusion unless
the entire required configuration space has been covered.


\subsection{Organization and the role of computation}

We first prove the finite trace and inertia inequality, then construct the
common local witness. A separate multiplicity calculation retains the
simple Gram block while grouping all other locations by multiplicity.
Only after those arguments are complete do we introduce the two kernels
and invoke pair correlation. This order exposes the interface between
finite linear algebra, local numerical certification and analytic number
theory.

The numerical part has two logically separate jobs. Rigorous interval
enclosures certify the local inequalities on the entire nonnegative gap
orthant; exact rational interval arithmetic then encloses the constants
in \eqref{eq:headlineintersection}--\eqref{eq:higherlocations}. The latter
calculation does not replace the former. Conversely, a local certificate
has no zero-count meaning until its matching kernel has passed through the
analytic limit. The supplement explains the full acceptance mathematics,
the finite coverage rules and the dependence on rigorous elementary-function
arithmetic.

\section{A finite trace inequality with a free clipping threshold}
\label{sec:finite}

\subsection{The operator attached to a conjugation-invariant multiset}

Let $\eta$ be a real even square-integrable function of compact support,
with $\int_\R\eta(s)^2\dd s=1$. Put $f=\eta^2$ and
\begin{equation}
 K(z)=\int_\R f(s)e^{-2\pi izs}\dd s,\qquad
 v_z(s)=\eta(s)e^{-2\pi izs}.\label{eq:kernelvectors}
\end{equation}
Compact support makes $K$ entire. It is even, real on the real axis, and
$K(0)=1$. We take the inner product to be linear in its first variable,
so that
\[
 \langle v_z,v_w\rangle=K(z-\overline w).
\]
The rank-one notation $v\otimes w^*$ denotes the operator
$h\mapsto\langle h,w\rangle v$.

Let $\mathcal Z$ be a finite conjugation-invariant multiset in $\C$, with
total multiplicity $N$. A location and each of its repeated copies are
distinguished when forming the following sum:
\begin{equation}
 A=\sum_{z\in\mathcal Z}v_z\otimes v_{\bar z}^{*}.
 \label{eq:A}
\end{equation}
Conjugation invariance gives $A=A^*$, and every summand has trace one.
All calculations may be made in the finite-dimensional span of the vectors
that occur, without any assertion of their linear independence. In
particular,
\begin{equation}
 \Tr A=N,\qquad
 H:=\|A\|_\HS^2=\sum_{z,w\in\mathcal Z}K(z-w)^2.
 \label{eq:traceenergy}
\end{equation}
For the second identity, if $A_z=v_z\otimes v_{\bar z}^*$, then
$\Tr(A_zA_w)=K(w-z)K(z-w)=K(z-w)^2$. Summation and selfadjointness
give the identity. The square in this formula is a complex square, not
the absolute square of each summand. Individual summands need not be
nonnegative; the total is a squared Hilbert--Schmidt norm.

Let $x_1\le\cdots\le x_n$ be the simple real locations in $\mathcal Z$.
Define
\begin{equation}
 P=\sum_{j=1}^n v_{x_j}\otimes v_{x_j}^*,\qquad
 M=(K(x_i-x_j))_{1\le i,j\le n},\qquad Q=A-P.
 \label{eq:simpleblock}
\end{equation}
The Gram matrix $M$ is positive semidefinite and has diagonal one. Its
order is $n$, even if its rank is smaller. The nonzero spectra of $M$
and $P$ agree, and $\Tr P=n$.

\subsection{Exceptional orbits and positive inertia}

Let $r_R$ count the multiple real locations, of respective multiplicities
$k_j\ge2$, and let $r_C$ count the unordered nonreal conjugate pairs,
of respective multiplicities $\ell_j\ge1$ at each member of a pair.
Set $b=r_R+r_C$. Each real multiple contributes one positive rank-one
term to $Q$. For a nonreal pair write
\[
 g_z=(v_z+v_{\bar z})/2,\qquad
 h_z=(v_z-v_{\bar z})/(2i).
\]
Its contribution is
\[
 \ell_z(v_z\otimes v_{\bar z}^*+v_{\bar z}\otimes v_z^*)
 =2\ell_z(g_z\otimes g_z^*-h_z\otimes h_z^*).
\]
Consequently $Q=Q_a-Q_b$ with $Q_a,Q_b\ge0$ and $\rank Q_a\le b$.
On $\ker Q_a$ the quadratic form of $Q$ is nonpositive, so the number
of its positive eigenvalues satisfies
\begin{equation}
 n_+(Q)\le b.\label{eq:inertia}
\end{equation}
These nonorthogonal orbit terms are used only to prove this inertia
bound. The square completion below instead uses the spectral positive
and negative parts of $Q$, which are orthogonal.

For $t\ge0$ define the clipped Gram defect
\begin{equation}
 D_t(M)=\Tr\bigl((M-I)^2-(M-(1+t)I)_+^2\bigr).
 \label{eq:defect}
\end{equation}
Here the positive part is defined spectrally. Expanding the first square
and retaining the zero eigenvalues of $M$ gives the exact identity
\begin{equation}
 D_t(M)=\Tr P^2-n-\Tr(P-(1+t)I)_+^2.\label{eq:defectP}
\end{equation}
The value is nonnegative: for an eigenvalue $x\ge0$ of $M$, the scalar
function $(x-1)^2-(x-1-t)_+^2$ is nonnegative.

For every $\Gamma\ge2$, with $t=\Gamma/2-1$, we now obtain
\begin{equation}
 H\ge\Gamma N-(\Gamma-1)n-\frac{\Gamma^2}{4}b+D_t(M).
 \label{eq:finiteinertia}
\end{equation}
Indeed, write $Q=Q_+-Q_-$ for its spectral decomposition. Subtracting
$\Gamma\Tr A$ from $\Tr A^2$ gives
\begin{align*}
 H-\Gamma N={}&\Tr P^2-\Gamma n
 +\Tr(Q_+^2-\Gamma Q_+)
 +\Tr(Q_-^2+\Gamma Q_--2PQ_-)
 +2\Tr(PQ_+).
\end{align*}
The last term is nonnegative. The first quadratic expression in $Q_+$
is at least $-\Gamma^2 b/4$, by \eqref{eq:inertia} and scalar square
completion. For a selfadjoint $X$ and $B\ge0$,
\begin{equation}
 \Tr(B^2-2XB)\ge-\Tr X_+^2.\label{eq:positivecompletion}
\end{equation}
To see this, diagonalize $X$: off-diagonal entries of $B$ add a
nonnegative amount to $\Tr B^2$, and each nonnegative diagonal entry
minimizes its scalar quadratic at the positive part of the corresponding
eigenvalue of $X$. Apply this with $X=P-\Gamma I/2$. Equation
\eqref{eq:defectP} completes the proof of \eqref{eq:finiteinertia}.

\subsection{The excess that the different counts measure}

Write
\begin{equation}
 u=N-n,\qquad e=u-2b
 =\sum_{j=1}^{r_R}(k_j-2)+2\sum_{j=1}^{r_C}(\ell_j-1)\ge0.
 \label{eq:excess}
\end{equation}
Use the definitions of $U,\Nd,M_{\le r},D_{\le r}$ from the
introduction, with ``critical'' replaced for now by ``real.'' The mass
excluded from the union consists precisely of the nonreal pairs with
$\ell_j\ge2$. Therefore
\begin{equation}
 N-U=2\sum_{\ell_j\ge2}\ell_j\le2e.\label{eq:unionexcess}
\end{equation}
For $r\ge2$, writing $L_{>r}=N-M_{\le r}$, we similarly have
\begin{equation}
 L_{>r}\le\frac{r+1}{r-1}e.\label{eq:massexcess}
\end{equation}
For a real location of multiplicity $k\ge r+1$, this follows from
$k/(k-2)\le(r+1)/(r-1)$. A nonreal pair of multiplicity
$\ell\ge r+1$ contributes $2\ell$ to the mass and $2(\ell-1)$ to
$e$, with an even smaller ratio. Summing proves the assertion.

The distinct count has a separate exact identity:
\begin{equation}
 N-\Nd=e+r_R,\qquad
 \Nd\ge\frac{N+n-e}{2}.\label{eq:distinctidentity}
\end{equation}
The second assertion uses $r_R\le b=(N-n-e)/2$. In particular it is
not an inference from a bound for $n$ alone: the excess term must be
retained until the trace inequality is substituted.

Suppose, for some $a,d,B\ge0$, a local argument has established
\begin{equation}
 D_t(M)\ge an-dS-B,\qquad
 S=x_n-x_1\quad(n\ge2),\quad S=0\quad(n\le1).
 \label{eq:affinedefect}
\end{equation}
Substituting $b=(u-e)/2$ into \eqref{eq:finiteinertia} yields
\begin{equation}
 H\ge(1+a)N-dS-B+
 \left(\Gamma-1-\frac{\Gamma^2}{8}-a\right)u
 +\frac{\Gamma^2}{8}e.\label{eq:masterexcess}
\end{equation}
At $\Gamma=4$ this is equivalently
\begin{equation}
 (1-a)n\ge2N-H-dS-B+2e.\label{eq:intersectionfinite}
\end{equation}
For the union, it is useful instead to choose
\begin{equation}
 \Gamma=4+\sqrt{8(1-a)},\qquad 0\le a<1,
 \label{eq:unionthreshold}
\end{equation}
which makes the coefficient of $u$ exactly zero. This algebraic choice,
and not a numerical sign test near zero, will be used in the limit.

\section{One witness for all overlapping windows}
\label{sec:witness}

\subsection{The variational formula}

Let $G\ge0$ be an $n\times n$ Hermitian matrix with diagonal one, and
write $Y=G-I$. For $t\ge1$,
\begin{equation}
 D_t(G)=\max_{T=T^*,\ \|T\|_\op\le t}
       \{2\Tr(YT)-\Tr(T^2)\}.\label{eq:dual}
\end{equation}
After diagonalizing $Y$, deleting the off-diagonal entries of $T$ does
not change the linear term, reduces its squared norm, and preserves the
constraint on its diagonal. The scalar maximum is
$y^2-(y-t)_+^2$ because $y\ge-1\ge-t$. This proves the formula.
The restriction $t\ge1$ is essential here: a symmetric norm ball would
also clip some negative eigenvalues if $t<1$.

Fix $m\ge2$ and nonnegative symmetric coefficients $c_{ij}$, with zero
diagonal, such that
\begin{equation}
 \sum_{i=1}^{m-h}c_{i,i+h}\le2\quad(1\le h<m),\qquad
 C_{\max}:=\max_i\sum_{j\ne i}c_{ij}>0.
 \label{eq:capacity}
\end{equation}
For each full consecutive block $B$ of length $m$, use local indices
$1,\ldots,m$ and define
\begin{gather}
 V_B=\sum_{i<j}c_{ij}|G_{B_iB_j}|^2,\qquad
 (Z_B)_{ij}=\tfrac m2c_{ij}G_{B_iB_j}\ (i\ne j),
 \quad (Z_B)_{ii}=0,\label{eq:localwitness}\\
 s_B=\max_i\sum_j|(Z_B)_{ij}|,\qquad
 \nu_B=\min(1,t/s_B),\qquad T_B=\nu_B Z_B,
 \label{eq:nu}
\end{gather}
where $\nu_B=1$ when $s_B=0$. The row-sum bound gives
$\|T_B\|_\op\le t$.

There are $m$ partitions of the integer line into consecutive $m$-blocks,
one for each residue of the first index modulo $m$. Restrict them to
$\{1,\ldots,n\}$, placing $T_B$ on every full block and zero on each
incomplete endpoint block. The resulting $m$ block-diagonal matrices
all have norm at most $t$. Their average $\overline T$ is therefore
admissible in \eqref{eq:dual}. Every full consecutive window occurs
exactly once in this average.

The linear term is
\[
 2\Tr(Y\overline T)=2\sum_B\nu_B V_B.
\]
For a fixed global pair of indices of separation $h$, its local
coefficients among the windows containing it form a subcollection of
$c_{1,1+h},\ldots,c_{m-h,m}$. Their sum is at most two. Weighted
Cauchy--Schwarz consequently gives
\[
 \left(\sum_{B\ni i,j}c_{ij}^{B}\nu_B\right)^2
 \le2\sum_{B\ni i,j}c_{ij}^{B}\nu_B^2.
\]
Since the corresponding entry of $\overline T$ is
$\frac12G_{ij}\sum_{B\ni i,j}c_{ij}^{B}\nu_B$, summing its absolute
square over both orientations of each pair yields
\[
 \Tr\overline T^2\le\sum_B\nu_B^2V_B.
\]
Thus one admissible witness proves the simultaneous estimate
\begin{equation}
 D_t(G)\ge\sum_B J_t(B),\qquad
 J_t(B):=(2\nu_B-\nu_B^2)V_B.\label{eq:commonwitness}
\end{equation}
This is where the overlap is paid for. No sum of separate spectral
pinchings, and no subsequent outer-block envelope, is required.

\subsection{The threshold-dependent cap and pressure transfer}

The row sum in \eqref{eq:nu} is controlled by the very energy that is
being retained. For each row, weighted Cauchy--Schwarz gives
\[
 s_B^2\le\frac{m^2}{4}C_{\max}V_B.
\]
Put
\begin{equation}
 b_m=\frac4{m^2C_{\max}}.\label{eq:bm}
\end{equation}
If $s_B\le t$, then $J_t(B)=V_B$. If $s_B\ge t$, then
\begin{equation}
 J_t(B)=\left(\frac{2t}{s_B}-\frac{t^2}{s_B^2}\right)V_B
 \ge b_m(2ts_B-t^2)\ge b_m t^2.\label{eq:badcap}
\end{equation}

Now take $G_{ij}=K(x_i-x_j)$ for real ordered points, and let
$g_j=x_{j+1}-x_j\ge0$. Suppose the local kernel inequality
\begin{equation}
 \sum_{i<j}c_{ij}\left|K(g_i+\cdots+g_{j-1})\right|^2
       +\sum_{j=1}^{m-1}p_jg_j\ge\epsilon
 \quad(g_j\ge0)\label{eq:rawrow}
\end{equation}
holds with $p_j\ge0$ and
\begin{equation}
 0\le\epsilon\le b_m t^2.\label{eq:capcondition}
\end{equation}
On an unclipped block, \eqref{eq:rawrow} bounds $J_t$ directly; on a
clipped block, \eqref{eq:badcap} and nonnegative pressure do so.
Sum over all $n-m+1$ full windows when $n\ge m$. Each elementary
gap is charged by at most one copy of each $p_j$, so its total
coefficient is at most $P_*:=\sum_{j=1}^{m-1}p_j$. Equation
\eqref{eq:commonwitness} then proves
\begin{equation}
 D_t(G)\ge\epsilon n-P_*S-(m-1)\epsilon.
 \label{eq:pressuretransfer}
\end{equation}
For $n<m$, the right side is nonpositive and the assertion follows
from $D_t\ge0$. This also treats empty and one-point Gram matrices.
The endpoint term is retained until the eventual height limit.

Only the unclipped domain $s_B\le t$ needs \eqref{eq:rawrow}; a
whole-orthant certificate is a convenient stronger premise. The
threshold condition nevertheless cannot be omitted. For the uniform
choice
\begin{equation}
 c_{i,i+h}=\frac2{m-h},\qquad
 C_{\max}=2\sum_{h=1}^{m-1}\frac1h,
 \label{eq:uniformweights}
\end{equation}
the endpoint row realizes the maximum. At $m=7$,
\begin{equation}
 b_7=\frac{40}{2401}.\label{eq:b7}
\end{equation}
Our cosine floor $a_0$ is smaller than $b_7$, whereas the two-mode
floor $a_1$ is larger. The latter is therefore not transferred at
$t=1$ by this argument. It is transferred at $t=2$, and at the
union threshold $t=1+\sqrt{2(1-a_1)}>2$.

\subsection{A bounded state correction}
\label{sec:potential}

The common witness also permits a bounded correction that telescopes
between neighboring windows. This observation is useful without
asserting that an advantageous correction has been found for any
particular unresolved numerical problem.

Put $q=m-1$, and let $\Phi:\R_{\ge0}^{q-1}\to\R$ have finite
oscillation $\Omega=\sup\Phi-\inf\Phi$. For $q=1$ its domain is a
singleton. Keep $t\ge1$, $p_j\ge0$ and $\epsilon\ge0$. Suppose
that for every nonnegative $q$-tuple,
\begin{equation}
 J_t(g)+\sum_{j=1}^{q}p_jg_j
 \ge\epsilon+\Phi(g_1,\ldots,g_{q-1})
                   -\Phi(g_2,\ldots,g_q).\label{eq:subaction}
\end{equation}
The state differences telescope once across the complete sequence of
windows. Consequently
\begin{equation}
 D_t(G)\ge\epsilon n-P_*S-q\epsilon-\Omega.
 \label{eq:potentialtransfer}
\end{equation}
For $n<m$ the right side is nonpositive and $D_t\ge0$ again suffices.
No regularity of $\Phi$ is needed. If the corresponding inequality
with $V$ in place of $J_t$ is known only for $s\le t$, the sufficient
condition $\epsilon+\Omega\le b_mt^2$ supplies the complementary
domain using \eqref{eq:badcap}.

There is also a useful uniform stability bound. Suppose two normalized
nonnegative densities have Fourier transforms differing by at most
$\rho$ on the real line. Both transforms have absolute value at most
one, so \eqref{eq:capacity} gives $|V-V'|\le4q\rho$ and
$|s-s'|\le mC_{\max}\rho/2$. The function
$s\mapsto\min(1,t/s)$, continuously extended at zero, is
$1/t$-Lipschitz. Since $V'\le2q$, it follows that
\begin{equation}
 |J_t-J_t'|\le
 \left(4q+\frac{2qmC_{\max}}t\right)\rho.\label{eq:Jstability}
\end{equation}
Thus a proved joined inequality can be smoothed with a controlled
floor loss. One must smooth that full inequality, not assume that the
boundary $s=t$ remains fixed. A dilation by $\sigma>0$ replaces $\Phi$ by
$\Phi(\sigma\,\cdot)$ and leaves its oscillation unchanged.

\section{Multiplicity accounting with the simple Gram block retained}
\label{sec:multiplicity}

Let $x_m$ count real locations of multiplicity $m$ and let $y_\ell$
count unordered nonreal pairs of multiplicity $\ell$ at each member.
Then $n=x_1$, $N=\sum_m m x_m+2\sum_\ell\ell y_\ell$, and
$\Nd=\sum_m x_m+2\sum_\ell y_\ell$. For $c\ge1$ define
\begin{equation}
 h_c(x)=c^2-(c-x)_+^2\quad(x\ge0).\label{eq:hc}
\end{equation}
This is nondecreasing and concave, and $h_c(0)=0$.

Include \emph{all} real locations in the positive operator
$P_R=\sum_{z\in\mathcal Z\cap\R}v_z\otimes v_z^*$, with
multiplicity. The positive inertia of $A-P_R$ is at most
$y:=\sum_\ell y_\ell$. Repeating the square completion in
Section~\ref{sec:finite}, now with $\Gamma=2c$, gives
\begin{equation}
 H\ge2cN-\Tr h_c(P_R)-c^2y.\label{eq:premaster}
\end{equation}
For clarity, the contribution of an eigenvalue $x$ of $P_R$ is
$x^2-2cx-(x-c)_+^2=-h_c(x)$, which proves the displayed form.

Use one weighted Gram column $\sqrt m\,v_z$ per distinct real
location of multiplicity $m$. Pinch this Gram matrix into the entire
simple-real block $M$ and one scalar block $[m]$ per multiple-real
location. Scalar concavity gives
\begin{equation}
 \Tr h_c(P_R)\le\Tr h_c(M)+\sum_{m\ge2}h_c(m)x_m.
 \label{eq:pinching}
\end{equation}
Indeed, diagonalize each compressed block and use its eigenvectors as
an orthonormal basis of the whole Gram space. Jensen's inequality
$h_c(\langle Gv,v\rangle)\ge\langle h_c(G)v,v\rangle$ follows
from the spectral measure of $G$ at $v$. Summing proves
\eqref{eq:pinching}. Operator concavity is not required. Gram zero
eigenvalues cause no discrepancy because $h_c(0)=0$.

Since $\Tr M=n$, the exact simple-block identity is
\[
 \Tr h_c(M)=(2c-1)n-D_{c-1}(M).
\]
Combining this with \eqref{eq:premaster} and \eqref{eq:pinching}
produces the corrected multiplicity inequality
\begin{equation}
 H\ge2cN-\sum_{m\ge1}h_c(m)x_m-c^2\sum_{\ell\ge1}y_\ell
           +D_{c-1}(M).\label{eq:multiplicitymaster}
\end{equation}
The scalar master without the final Gram term is the comparison
framework in \cite{Waddle}. The retained term is indispensable for
the improvements made here.

\subsection{Mass of multiplicity at most two}

Take $c=3$, and suppose $D_2(M)\ge an-dS-B$ with $0\le a\le1/2$.
In \eqref{eq:multiplicitymaster}, the real-location debits are
$5-a$ for $m=1$, $8$ for $m=2$, and $9$ for $m\ge3$; the debit
per nonreal pair is $9$. Per unit mass of multiplicity at most two,
these are at most $5-a$: the four possible values are
$5-a,4,9/2,9/4$. For all higher multiplicities the debit per unit
mass is at most $3$. Hence
\[
 H+dS+B\ge6N-(5-a)M_{\le2}-3(N-M_{\le2}),
\]
and therefore
\begin{equation}
 M_{\le2}\ge\frac{3N-H-dS-B}{2-a}.\label{eq:mass2finite}
\end{equation}
This is a mass estimate; it is not yet a proportion among distinct
locations.

\subsection{Costs per location and the distinct denominator}

Assume a finite upper budget
\begin{equation}
 H+dS+B\le\mathcal T N,\qquad
 \alpha=2c-\mathcal T,\quad \beta=\mathcal T-1-a.\label{eq:upperbudget}
\end{equation}
With $D_{c-1}(M)\ge an-dS-B$, equation
\eqref{eq:multiplicitymaster} implies
\begin{equation}
 \sum_{m\ge1}\lambda_m x_m+\sum_{\ell\ge1}\omega_\ell y_\ell\le0,
 \quad
 \lambda_m=\alpha m-h_c(m)+a\ind_{m=1},\quad
 \omega_\ell=2\alpha\ell-c^2.\label{eq:locationcosts}
\end{equation}
Suppose that for a fixed cutoff $r$ and some $\gamma_r>0$, these
costs are at least $-\beta$ per low-multiplicity location and at
least $\gamma_r-\beta$ per high-multiplicity location. A nonreal
pair contains two locations, so its required bounds are twice the
corresponding quantities. Then \eqref{eq:locationcosts} gives
$0\ge-\beta\Nd+\gamma_r(\Nd-D_{\le r})$, whence
\begin{equation}
 \frac{D_{\le r}}{\Nd}\ge1-\frac\beta{\gamma_r}.
 \label{eq:locationratio}
\end{equation}
This deduction requires $\Nd>0$, as will hold at sufficiently large
heights.

For $r=2$ and $c=3$, sufficient conditions are
\begin{equation}
 \beta\ge0,\quad \alpha=6-\mathcal T>0,\quad
 3-\mathcal T-a\ge0,\quad1-2a\ge0,\quad
 \gamma_2=8-2\mathcal T-a>0.\label{eq:guards2}
\end{equation}
Here $\lambda_1=-\beta$,
$\lambda_2+\beta=3-\mathcal T-a$, and for $m\ge3$
\[
 \lambda_m+\beta-\gamma_2=\alpha(m-3)\ge0.
\]
For pairs, $\omega_1+2\beta=1-2a\ge0$, and increasing
$\ell$ only increases the cost. At the first high multiplicity,
$\omega_3+2\beta-2\gamma_2=9$, which remains positive for
$\ell\ge3$. These facts check every multiplicity, not a finite
selection of them, and give
\begin{equation}
 \frac{D_{\le2}}{\Nd}\ge
 1-\frac{\mathcal T-1-a}{8-2\mathcal T-a}.\label{eq:distinct2finite}
\end{equation}

For $r\ge3$, let $3<c<4$, and assume
\begin{gather}
 \alpha>0,\quad\beta\ge0,\quad
 3-\mathcal T-a\ge0,\quad8-2\mathcal T-a\ge0,\label{eq:guardsr1}\\
 4\alpha+\beta-c^2\ge0,\qquad
 q_c:=2c-1-c^2/2-a\ge0,\qquad
 \gamma_r:=(r+1)\alpha+\beta-c^2>0.\label{eq:guardsr2}
\end{gather}
For $m=1,2,3$, the values of $\lambda_m+\beta$ are respectively
$0$, $3-\mathcal T-a$, and $8-2\mathcal T-a$. For $m\ge4$,
$\lambda_m+\beta=\alpha m+\beta-c^2$, which is nonnegative
at $m=4$ and increasing. At $m=r+1$ it equals $\gamma_r$.
For a nonreal pair the smallest value of $\omega_\ell+2\beta$
is $2q_c\ge0$. At $\ell=r+1$ it exceeds $2\gamma_r$ by
exactly $c^2$. Thus \eqref{eq:locationratio} holds for every fixed
$r\ge3$ under the displayed conditions.

At $2c=\Gamma$ from \eqref{eq:unionthreshold}, write
$Z=\Gamma^2/8$ and $E=\beta/Z$. The identity
$\Gamma-1-Z-a=0$ gives $\alpha=Z-\beta$ and
\begin{equation}
 \gamma_r=Z(r-1-rE).\label{eq:gammar}
\end{equation}
This explains the denominator in \eqref{eq:higherlocations}. It is
obtained by per-location accounting, not by dividing a mass estimate
by a distinct count.

\section{The two kernels and their certified local inequalities}
\label{sec:kernels}

\subsection{Normalization and the exact energy cost}

On $[-1/2,1/2]$, define
\begin{equation}
 f_0(s)=\frac{\cos(\sqrt2s)}{\sqrt2\sin z},\qquad
 f_1(s)=f_0(s)+u_1\cos(2\pi s)+u_2\cos(4\pi s),
 \qquad z=1/\sqrt2,\label{eq:densities}
\end{equation}
and extend both functions by zero outside that interval. The coefficients
$u_1,u_2$ are those in \eqref{eq:R0}. Direct integration gives
$\int f_0=\int f_1=1$. These are nonnegative densities, not signed
test functions: on the interval $\cos(\sqrt2s)\ge3/4$ and
$\sqrt2\sin z\le1$, so $f_0\ge3/4$ and
$f_1\ge3/4-|u_1|-|u_2|=71/100>0$.

For a real even compactly supported density $f\in L^1\cap L^2$, define
\begin{equation}
 \mathcal R(f)=\int_\R f(s)^2\dd s+
 \iint_{\R^2}|s-t|f(s)f(t)\dd s\dd t.\label{eq:energyfunctional}
\end{equation}
The energy of $f_0$ is $R_0$ in \eqref{eq:R0}. A short calculation
also explains why its perturbation cost is exactly quadratic. On the
support interval put $W_0(s)=\int|s-t|f_0(t)\dd t$. Then
$W_0''=2f_0$ and $f_0''=-2f_0$. Evenness shows that $f_0+W_0$
is constant. Its value at $s=1/2$ is
$z\cot z+1/2=R_0$, so
\[
 \mathcal R(f_0)=\int f_0(f_0+W_0)=R_0.
\]
If $h$ has integral zero on the interval, the linear variation of
$\mathcal R(f_0+h)$ vanishes. Write
$H_h(s)=\int_{-1/2}^s h(t)\dd t$, so $H_h$ vanishes at both
endpoints. With $W_h(s)=\int|s-t|h(t)\dd t$ one has
$W_h'=2H_h$ and hence, by integration by parts,
\[
 \iint |s-t|h(s)h(t)\dd s\dd t=-2\int H_h(s)^2\dd s.
\]
For $h(s)=\sum_j u_j\cos(2\pi js)$, trigonometric orthogonality
therefore gives
\begin{equation}
 \mathcal R(f_0+h)-R_0
 =\int h^2-2\int H_h^2
 =\sum_j\left(\frac12-\frac1{4\pi^2j^2}\right)u_j^2.
 \label{eq:energyperturbation}
\end{equation}
The functional is quadratic, so there is no omitted higher-order
remainder. In particular $\mathcal R(f_1)=R_1=R_0+\Delta$.

The two pieces of the energy must be kept separately when the support
is dilated. Put $A_j=\int f_j^2$ and $B_j=R_j-A_j$ for $j=0,1$.
Direct integration yields
\begin{align}
 A_0&=\frac1{4\sin^2z}+\frac{z\cot z}{2},\label{eq:A0}\\
 A_1&=A_0+2\sum_{j=1}^2
       \frac{(-1)^{j+1}u_j}{2\pi^2j^2-1}
             +\frac12\sum_{j=1}^2u_j^2.\label{eq:A1}
\end{align}
For example, the cross term follows from
$\int f_0(s)\cos(2\pi js)\dd s=(-1)^{j+1}/(2\pi^2j^2-1)$.
If $f_{j,\sigma}(s)=\sigma^{-1}f_j(s/\sigma)$, then
\begin{equation}
 \mathcal R(f_{j,\sigma})=R_j(\sigma)
 :=\frac{A_j}{\sigma}+B_j\sigma.\label{eq:dilatedenergy}
\end{equation}
Both coefficients are positive. The use of $A_1,B_1$, rather than
the cosine coefficients $A_0,B_0$, is necessary for the perturbed
short-height calculation.

\subsection{The local inputs}

Let $k_j=\widehat f_j$ with the Fourier convention
\eqref{eq:kernelvectors}. Writing $\operatorname{sinc}(x)=\sin(x)/x$
with value one at zero, their entire formulas are
\begin{align}
 k_0(x)&=\frac{\operatorname{sinc}(\pi x-z)
                    +\operatorname{sinc}(\pi x+z)}{2\sqrt2\sin z},
 \label{eq:k0sinc}\\
 k_1(x)&=k_0(x)+\frac12\sum_{j=1}^2u_j
       \bigl\{\operatorname{sinc}(\pi(x-j))
                     +\operatorname{sinc}(\pi(x+j))\bigr\}.
 \label{eq:k1sinc}
\end{align}
Define the seven-point energy for six nonnegative gaps by
\begin{equation}
 V_{7,j}(g)=\sum_{h=1}^{6}\frac2{7-h}
     \sum_{i=1}^{7-h}k_j(g_i+\cdots+g_{i+h-1})^2.
 \label{eq:sevenenergy}
\end{equation}
The two computer-assisted local inequalities are
\begin{align}
 V_{7,0}(g)+\frac1{3000}\sum_{i=1}^6g_i
 &\ge\frac{19}{5000},\label{eq:cosinerow}\\
 V_{7,1}(g)+\frac1{500}\sum_{i=1}^6g_i
 &\ge\frac7{400}.\label{eq:twomoderow}
\end{align}
Both hold for every $g\in\R_{\ge0}^6$, including zero gaps and
unbounded configurations. Equation~\eqref{eq:cosinerow} is Ainta's
certificate \cite{Ainta}. Equation~\eqref{eq:twomoderow} is the
two-mode certificate used for the present union and multiplicity bounds.
The supplement gives the rigorous finite verification in detail.

Here is the structure of that verification, to clarify what is asserted.
Nonnegative pressure first reduces a possible strict counterexample to a
compact simplex. A single-coordinate lower test cuts this down further.
Outward scalar enclosures on closed grid cells give lower bounds for
every squared-kernel term. A sum of $h$ gaps requires all $h$ final
scalar cells; the extra $h-1$ indices at a summed upper endpoint are
included. Adaptive binary subdivision covers the resulting boxes.
Interval energy, whole-box Taylor bounds and rigorous positive-definiteness
tests each supply sufficient acceptance rules. Removable singularities in
\eqref{eq:k0sinc}--\eqref{eq:k1sinc} use convergent power series with
explicit remainders. An unaccepted terminal box, a time limit or a failed
arithmetic guard gives no certificate. The recorded completed covers
exhaust their domains under these rules. This is stronger than checking
the objective at grid points, and the proof does not treat a digest or a
node count as a substitute for coverage.

In \eqref{eq:pressuretransfer}, these rows give $a=a_j$ and
$d=6p=d_j$, with endpoint debit $B=6a_j$. The first may be used at
$t=1$. For the second,
\[
 \frac{a_1}{b_7}=\frac{16807}{16000}>1,
 \qquad a_1<4b_7<b_7(1+\sqrt{2(1-a_1)})^2.
\]
Thus its admissible thresholds include the two required in
Sections~\ref{sec:finite} and \ref{sec:multiplicity}. These are
separate uses of its \emph{same} local certificate, not sums of
corrections from different densities.

\section{From the finite operator to zeta zeros}
\label{sec:analytic}

\subsection{The fixed-test pair-correlation input}

We state precisely the analytic theorem being used. If $g$ is a fixed
real even integrable function supported in $[-1,1]$ and Lipschitz at
zero, the test-function form of \cite[Lemma 5]{BGST}, also stated in
\cite[Lemma 3.1]{Lamzouri}, is, with $L=\log X$,
\begin{align}
 &\sum_{0<\gamma,\gamma'\le X}
 \widehat g\!\left(\frac{i(\rho-\rho')L}{2\pi}\right)
       \frac4{4-(\rho-\rho')^2}\notag\\
 &\hspace{12mm}=\frac{XL}{2\pi}
  \left\{g(0)+\int_\R|\alpha|g(\alpha)\dd\alpha
                      +O_g(L^{-1/2})\right\}.
 \label{eq:BGST}
\end{align}
The zeros are counted with multiplicity, and the formula is unconditional.
We use it only for fixed smooth functions supported strictly inside
$(-1,1)$. No new proof of \eqref{eq:BGST} is claimed here.

Take a real even smooth $\eta$ supported strictly in
$(-\sigma/2,\sigma/2)$, with $0<\sigma<1$ and $\int\eta^2=1$.
Set $f=\eta^2$, $K=\widehat f$ and $G=f*f$. Both $G$ and $G''$
are admissible fixed tests. For
$w=i(\rho-\rho')L/(2\pi)$, integration by parts gives
$\widehat{G''}(w)=-4\pi^2w^2K(w)^2$. Consequently the exact
identity
\begin{equation}
 K(w)^2=\left(\widehat G(w)-\frac{\widehat{G''}(w)}{4L^2}\right)
                 \frac4{4-(\rho-\rho')^2}\label{eq:weightremoval}
\end{equation}
removes the rational weight. Its denominator is nonzero: the real
part of the difference of two nontrivial zeros lies in $(-1,1)$,
so that difference cannot be $2$ or $-2$.

Apply \eqref{eq:BGST} \emph{separately} to $G$ and $G''$, then
combine their sums by \eqref{eq:weightremoval}. The second term is
$O_f(X/L)$, while the first has leading coefficient
\[
 G(0)+\int|\alpha|G(\alpha)\dd\alpha=\mathcal R(f).
\]
Evenness of $f$ changes the double integral with $|s+t|$ into that
with $|s-t|$. Thus, for the multiset
\begin{equation}
 \mathcal Z_X=\left\{\frac{i(\rho-1/2)\log X}{2\pi}:
                             0<\gamma\le X\right\},
 \label{eq:zetamultiset}
\end{equation}
equation~\eqref{eq:traceenergy} satisfies
\begin{equation}
 \frac{H_f(X)}{N(X)}\longrightarrow\mathcal R(f).
 \label{eq:energylimit}
\end{equation}
The functional equation and conjugation of $\zeta$ preserve the
multiplicity under $\rho\mapsto1-\overline\rho$. That map acts
as conjugation on \eqref{eq:zetamultiset}; its real elements are
exactly the critical-line zeros. Thus all finite counts above have
the interpretations \eqref{eq:counts1}--\eqref{eq:counts3}.

For completeness, the Riemann--von Mangoldt formula gives
\begin{equation}
 N(X)=\frac X{2\pi}\log\frac X{2\pi}-\frac X{2\pi}+O(\log X)
     =\frac{X\log X}{2\pi}+O(X),\label{eq:RVM}
\end{equation}
with the usual endpoint convention; see \cite{Titchmarsh} or
\cite[Lemma 2]{BGST}. The simple-real span obeys
$S\le X\log X/(2\pi)$, including when it is defined to be zero.
Therefore
\begin{equation}
 \limsup_{X\to\infty}S/N(X)\le1.\label{eq:spanlimit}
\end{equation}
Only an upper bound is needed, since every span debit is nonnegative.

\subsection{Smoothing without changing the test with height}
\label{sec:smoothing}

The densities \eqref{eq:densities} have sharp endpoints and cannot
simply be inserted as smooth strict-support windows. We now justify
their use in the final constants. Fix $j\in\{0,1\}$, a dilation
$0<\sigma<1$ and a slack $0<\delta<1$. Let
$F=f_{j,\sigma}$ and let $\psi$ be a nonnegative even smooth cutoff
supported strictly inside $(-\sigma/2,\sigma/2)$. Set
\begin{equation}
 f'=\frac{\psi^2F}{\int\psi^2F},\qquad
 \eta'=\psi\sqrt{F/\int\psi^2F}.\label{eq:smoothing}
\end{equation}
Choose cutoffs tending to one on every compact interior subinterval,
bounded between zero and one. Strict positivity and smoothness of $F$
on that interior make $\eta'$ smooth after extension by zero. The
densities $f'$ tend to $F$ in $L^1$ and $L^2$, and their energies
tend to $R_j(\sigma)$; the second part of
\eqref{eq:energyfunctional} is continuous in $L^1$ on a common
compact support.

If $\tau=\|f'-F\|_1$, their Fourier transforms differ by at most
$\tau$ on the \emph{real} axis. Both have absolute value at most
one there. The sum of all coefficients in a seven-point row is
$12$, so the squared-kernel energy changes by at most $24\tau$.
Choose the cutoff sufficiently close that $24\tau\le\delta a_j$.
Applying the original local row to the dilated gaps gives
\begin{equation}
 V_{7,f'}(g)+\sigma\frac{d_j}{6}\sum_{i=1}^6g_i
       \ge a'_j:=(1-\delta)a_j.\label{eq:smoothedrow}
\end{equation}
Only real transform arguments occur in a simple-real Gram block, so
no uniform comparison at complex zeta arguments is being assumed.

The common-witness transfer now reads
\begin{equation}
 D_t(M_{f'})\ge a'_j n-\sigma d_j S-6a'_j.
 \label{eq:smootheddefect}
\end{equation}
The pressure has become $\sigma d_j$; there is no division by
$1-\delta$. At $j=0$ use $t=1$. At $j=1$ use $t=2$ for
multiplicity two, and for the union recompute
\begin{equation}
 \Gamma'=4+\sqrt{8(1-a'_1)},\qquad
 t'=\Gamma'/2-1,\qquad Z'=\Gamma'^2/8.\label{eq:slackthreshold}
\end{equation}
Here $t'>2$ and the cap remains valid. The coefficient of $u$ in
\eqref{eq:masterexcess} is still exactly zero.

Every application of \eqref{eq:BGST} and \eqref{eq:energylimit}
has the smooth density, dilation, slack and threshold fixed before
$X\to\infty$. One then removes the cutoff at fixed slack, sends
$\delta\downarrow0$, and finally sends $\sigma\uparrow1$.
These successive limits of numerical lower bounds require no
uniform analytic error as the test changes. An arbitrarily small
fixed upper-budget margin for the distinct ratios is handled in
the same way below.

\section{Evaluation of the global proportions}
\label{sec:global}

\subsection{The simple-critical and total-distinct counts}

Use the smoothed cosine density in \eqref{eq:intersectionfinite},
with $a=a'_0$, $d=\sigma d_0$ and $B=6a'_0$. Since $e\ge0$,
the height limit, \eqref{eq:energylimit} and \eqref{eq:spanlimit}
give
\[
 \liminf_{X\to\infty}\frac{n(X)}{N(X)}
 \ge\frac{2-\mathcal R(f')-\sigma d_0}{1-a'_0}.
\]
Take the successive approximation limits described in
Section~\ref{sec:smoothing}; this proves
\eqref{eq:headlineintersection} with its exact constant $C$.

To obtain the distinct count, substitute the \emph{full} finite
inequality into \eqref{eq:distinctidentity} before dropping $e$:
\begin{align*}
 \Nd&\ge\frac12\left\{N+
    \frac{2N-H-\sigma d_0S-6a'_0}{1-a'_0}
            +\left(\frac2{1-a'_0}-1\right)e\right\}\\
 &\ge\frac12\left\{N+
    \frac{2N-H-\sigma d_0S-6a'_0}{1-a'_0}\right\}.
\end{align*}
The coefficient discarded is positive. The same limits now prove
\eqref{eq:headlinedistinct}.

\subsection{The union and higher-multiplicity mass}

Use the smoothed two-mode density and \eqref{eq:slackthreshold}.
Equation~\eqref{eq:masterexcess} gives
\[
 Z'e\le H+\sigma d_1S+6a'_1-(1+a'_1)N.
\]
For this fixed smooth test it follows that
\[
 \limsup_{X\to\infty}\frac eN
 \le\frac{\mathcal R(f')+\sigma d_1-1-a'_1}{Z'}.
\]
Equation~\eqref{eq:unionexcess} and the successive approximation
limits give \eqref{eq:headlineunion}. Equation~\eqref{eq:massexcess}
gives \eqref{eq:highermass} for each fixed $r\ge3$. Although
\eqref{eq:massexcess} is also valid at $r=2$, the direct estimate
\eqref{eq:mass2finite}, at $t=2$ with the same smoothed row, is
stronger there and yields \eqref{eq:headlinemass}.

\subsection{Distinct low-multiplicity locations}

A fixed finite budget avoids any uncontrolled division of an
$o(N)$ term by $\Nd$. Fix $\xi>0$ and, for a fixed smoothed
two-mode density, set
\[
 \mathcal T=\mathcal R(f')+\sigma d_1+\xi.
\]
Equations~\eqref{eq:energylimit} and \eqref{eq:spanlimit}, together
with $6a'_1/N\to0$, show that
$H+\sigma d_1S+6a'_1\le\mathcal T N$ for all sufficiently large
$X$. For sufficiently small $\xi$, slack, cutoff error and
$1-\sigma$, the strict versions of the guards
\eqref{eq:guards2} hold. Apply the finite ratio
\eqref{eq:distinct2finite} first, then pass to the height limit
and remove $\xi$ and the approximation parameters. This proves
\eqref{eq:headlinelocations}.

For each fixed $r\ge3$, use $c=\Gamma'/2$ and
\eqref{eq:guardsr1}--\eqref{eq:guardsr2}. The same procedure and
\eqref{eq:gammar} give \eqref{eq:higherlocations}. The limiting
guards can be checked without delicate numerical signs. For example,
\[
 \tfrac54\le R_0<\tfrac43,\qquad
 0<\Delta<\frac{109}{250000},\qquad 1<T_1<1.35,
 \qquad 6<\Gamma<7
\]
give $\beta>0$, $3-T_1-a_1>0$, $8-2T_1-a_1>0$ and
$\alpha=\Gamma-T_1>0$. Moreover, using $c^2=2Z$ and
$\alpha=Z-\beta$,
\[
 4\alpha+\beta-c^2=2Z-3\beta>0,
\]
since $Z>5$ and $\beta<0.35$. The remaining guard $q_c=0$ is
the exact identity \eqref{eq:unionthreshold}. Thus the strict
conditions persist under the small perturbations just used, with
this identity retained symbolically. The fine decimal enclosures
are independent of these coarse qualitative checks.

One can also see directly why the denominator cannot become too small
for a fixed test. Cauchy--Schwarz on the nonzero eigenvalues of $A$
and its rank-one decomposition by distinct locations give
\begin{equation}
 N^2\le\rank(A)H\le\Nd H.\label{eq:ranksafeguard}
\end{equation}
Since $H=O_f(N)$, this implies $\Nd\gg_f N$. We nevertheless
used a fixed upper budget so that the per-location inequalities
remain exact throughout the argument.

\subsection{Exact arithmetic and strict comparison signs}

The numerical constants in the introduction are enclosed by rational
intervals, not by floating-point rounding to a presumed number of
correct digits. The verifier constructs $\pi$ from Machin's identity
and alternating arctangent series, bounds square roots by rational
squaring, and evaluates sine and cosine with explicit Taylor
remainders. Interval division is used only after excluding zero
from the denominator. It evaluates \eqref{eq:R0}--\eqref{eq:E}
and the five exact expressions independently of the local cover
program. In particular it checks the cap inequalities and the
strict margins between each exact lower endpoint and its displayed
decimal. 

The signs of the improvements over the named basic formulas also
have a short analytic proof. Put
\[
 Z_0=3+2\sqrt2,\qquad
 \eta_*=a_1-d_1-\Delta,\qquad B=(R_0-1)/Z_0.
\]
Then $\beta=R_0-1-\eta_*$. Since $a_1<1/2$,
\[
 0<Z_0-Z=a_1\left(1+
          \frac{2\sqrt2}{1+\sqrt{1-a_1}}\right)<3a_1.
\]
Also $\eta_*>633/125000$, whereas
\[
 \frac{(R_0-1)(Z_0-Z)}{Z_0}<\frac{a_1}{5}
     =\frac7{2000}<\frac{633}{125000}.
\]
It follows that $E<B$, which proves the strict union improvement
and the stated higher-cutoff family comparisons. The estimates
$1<R_0<4/3$ used here follow from $\cos z\ge3/4$,
$\sin z/z\le1$, and, for the upper bound, the alternating-series
inequality
$\cos z<(5/6)(\sin z/z)$ at $z^2=1/2$.
For example $\cos z\le1-z^2/2+z^4/24=73/96$, while
$(5/6)(\sin z/z)\ge(5/6)(1-z^2/6)=55/72>73/96$.

The comparison with the mass-two formula $(3-R_0)/2$ has numerator
\[
 K_*=a_1(3-R_0)-2(\Delta+d_1)
 >\frac7{240}-2\left(\frac{109}{250000}+\frac3{250}\right)>0
\]
over the positive denominator $2(2-a_1)$. The comparison with the
distinct-location formula $1-(R_0-1)/(8-2R_0)$ has numerator
$3K_*$ over
$(8-2R_0)(8-2T_1-a_1)>0$. These identities explain the gain
without relying on a long decimal expansion. The exact scalar
calculation supplies the finer margins in
Section~\ref{subsec:literaturecomparison}.

\section{Short-height intervals with fixed support margin}
\label{sec:short}

Fix $0<\theta<1$ and let $H_X=X^\theta$. In this section all counts
are restricted to $X<\gamma\le X+H_X$, with their definitions
and denominators otherwise unchanged. We append a subscript $I$
when needed. For every fixed $0<\sigma<\theta$ and fixed smooth
real even test $g$ supported in $[-\sigma,\sigma]$, Wang's
Theorem~2.2 \cite{Wang} states
\begin{align}
 &\sum_{X<\gamma,\gamma'\le X+H_X}
 \widehat g\!\left(\frac{i(\rho-\rho')\log X}{2\pi}\right)
       \frac4{4-(\rho-\rho')^2}\notag\\
 &\hspace{8mm}=\frac{H_X\log X}{2\pi}
       \left(g(0)+\int|\alpha|g(\alpha)\dd\alpha\right)
                +O_g(H_X+X^\sigma\log^2X).
 \label{eq:Wang}
\end{align}
This is a separate analytic input, not a subtraction of two
full-height correlation asymptotics. Its normalized error tends
to zero for every fixed strict support margin:
\[
 \frac{H_X+X^\sigma\log^2X}{H_X\log X}
 =O\bigl((\log X)^{-1}+X^{\sigma-\theta}\log X\bigr)=o(1).
\]
Apply \eqref{eq:weightremoval} to the two fixed tests $G,G''$
as before. The difference form of \eqref{eq:RVM} gives
\[
 N_I(X)=\frac{H_X\log X}{2\pi}+O(H_X+\log X),
 \qquad S\le\frac{H_X\log X}{2\pi}.
\]
It follows that $H_f/N_I\to\mathcal R(f)$ and
$\limsup S/N_I\le1$. The finite argument is unchanged.

Define
\begin{equation}
 T_j(\theta)=\frac{A_j}{\theta}+(B_j+d_j)\theta\quad(j=0,1),
 \qquad E_\theta=\frac{T_1(\theta)-1-a_1}{Z}.
 \label{eq:shortconstants}
\end{equation}
Smooth first at a fixed $\sigma<\theta$, take the height limit
with every test fixed, remove the cutoff and floor slack, and only
then let $\sigma\uparrow\theta$. The resulting lower limiting
bounds are
\begin{align}
 \liminf\frac{n_I}{N_I}&\ge\frac{2-T_0(\theta)}{1-a_0},&
 \liminf\frac{\Nd_I}{N_I}&\ge
       \frac12\left(1+\frac{2-T_0(\theta)}{1-a_0}\right),
 \label{eq:shortintersection}\\
 \liminf\frac{U_I}{N_I}&\ge1-2E_\theta,&
 \liminf\frac{M_{\le2,I}}{N_I}&\ge\frac{3-T_1(\theta)}{2-a_1},
 \label{eq:shortunion}\\
 \liminf\frac{M_{\le r,I}}{N_I}&\ge
       1-\frac{r+1}{r-1}E_\theta\qquad(r\ge3\text{ fixed}).
 \label{eq:shorthigher}
\end{align}
All lower limits here are as $X\to\infty$. A negative right side
may be replaced by zero. These mass and total-distinct statements
hold for each fixed $0<\theta<1$ under the stated short-height
input, although at small $\theta$ some give no nontrivial lower
bound.

For distinct low-multiplicity locations, the parameter guards must
also survive this substitution. The bound
\begin{equation}
 \liminf\frac{D_{\le2,I}}{\Nd_I}\ge
 1-\frac{T_1(\theta)-1-a_1}{8-2T_1(\theta)-a_1}
 \label{eq:shortdistinct2}
\end{equation}
holds whenever the strict versions of the nonconstant guards in
\eqref{eq:guards2} hold with $a=a_1$ and
$\mathcal T=T_1(\theta)$. More explicitly, require
\[
 T_1(\theta)-1-a_1>0,\quad 6-T_1(\theta)>0,\quad
 3-T_1(\theta)-a_1>0,\quad8-2T_1(\theta)-a_1>0;
\]
the fixed condition $1-2a_1>0$ already holds.
For $r\ge3$, the analogous formula is
\begin{equation}
 \liminf\frac{D_{\le r,I}}{\Nd_I}\ge
           1-\frac{E_\theta}{r-1-rE_\theta},\label{eq:shortdistinctr}
\end{equation}
provided \eqref{eq:guardsr1}--\eqref{eq:guardsr2} hold with
$c=\Gamma/2$, $a=a_1$, $\mathcal T=T_1(\theta)$, strictly for
the nonconstant inequalities and with $q_c=0$ retained exactly.
These conditions hold for $\theta$ sufficiently close to one by
the global checks and continuity. We do not assert them for every
$\theta\in(0,1)$. The fixed upper-budget construction in
Section~\ref{sec:global} justifies the distinct denominator here
as well. None of these statements is uniform as $\theta$ varies
with height, and none sets $\sigma=\theta$ inside
\eqref{eq:Wang}.

\section{Scope of the accompanying local extension}
\label{sec:scope}

The finite transfer permits other window lengths, profiles and
distance-capacity weights, but its premise remains a uniform local
inequality. To make that boundary concrete, the supplement also
studies a thirteen-point objective with six explicit Fourier modes,
two gap pressures and target $31/5000$. Scalar exclusion partitions
the gap range of a hypothetical strict counterexample into short
and long components. Complete local certificates establish the
all-long region and every region with exactly one short gap,
including all ten interior positions. Their proof uses disjoint
triples and an exact convex quadratic lower bound on the remaining
small cells.

The resulting lower bounds are $0.006424625$ on the all-long
region, $0.00623$ on the two endpoint-one-short regions and
$0.0062$ on each interior-one-short region, on the precise domains
specified in the supplement. Thus a remaining strict counterexample
must have at least two short gaps. Whether such a configuration
exists remains unresolved in this work. The full thirteen-point inequality is not
asserted, and none of its prospective zero-count consequences is
included in \eqref{eq:headlineintersection}--\eqref{eq:higherlocations}.

The distinction is structural. The common-witness inequality can
transfer a proved local floor, and the bounded state correction can
transfer a proved subaction. Neither constructs a missing floor or
a missing potential. In the present application the two completed
seven-point inputs already suffice for every global bound stated
in the article. The larger-window result records an additional
proved local reduction without enlarging those conclusions.

\section*{Acknowledgments}
The author is indebted to his family for their extraordinary patience and support. He is grateful to Dr. C. Oommen and Prof. B. N. Raghunandan of the Department of Aerospace Engineering, Indian Institute of Science, for their mentorship and for fostering the intellectual latitude that enabled this independent research.

\section*{Note on AI Use}
The author declares the use of consumer-grade Gemini 3.1 and GPT-6 models for prior-art search, developing some ideas and proofs, computer-assisted verification, drafting and LaTeX typesetting; the author takes full and sole responsibility for all contents of this work.

\begin{thebibliography}{99}
\raggedright

\bibitem{AF} L. Alp\"oge and R. Furman,
\emph{More than two thirds of the zeta zeros are simple and on the critical line},
\href{https://arxiv.org/abs/2608.13637v2}{arXiv:2608.13637v2} (2026).

\bibitem{Ainta} GitHub contributor \texttt{ainta},
\emph{More than 67.3\% of the zeros of the Riemann zeta function
are simple and lie on the critical line}, public research repository,
\href{https://github.com/ainta/zeta-simple-zeros/tree/040c5e899e658aed7b56a2a87f501798fe10761d}{commit 040c5e8} (2026).

\bibitem{Aryan} F. Aryan,
\emph{On an extension of the Landau--Gonek formula},
J. Number Theory \textbf{233} (2022), 389--404,
\href{https://doi.org/10.1016/j.jnt.2021.06.015}{doi:10.1016/j.jnt.2021.06.015}.

\bibitem{BGST} S. A. C. Baluyot, D. A. Goldston,
A. I. Suriajaya and C. L. Turnage-Butterbaugh,
\emph{An unconditional Montgomery theorem for pair correlation of zeros
of the Riemann zeta-function}, Acta Arith. \textbf{214} (2024), 357--376,
\href{https://doi.org/10.4064/aa230612-20-3}{doi:10.4064/aa230612-20-3}.

\bibitem{BGSTI} S. A. C. Baluyot, D. A. Goldston,
A. I. Suriajaya and C. L. Turnage-Butterbaugh,
\emph{Pair Correlation of Zeros of the Riemann Zeta Function I:
Proportions of Simple Zeros and Critical Zeros},
\href{https://arxiv.org/abs/2501.14545v3}{arXiv:2501.14545v3} (2026).

\bibitem{CCLM} E. Carneiro, V. Chandee, F. Littmann and M. B. Milinovich,
\emph{Hilbert spaces and the pair correlation of zeros of the
Riemann zeta-function}, J. Reine Angew. Math. \textbf{725} (2017),
143--182,
\href{https://doi.org/10.1515/crelle-2014-0078}{doi:10.1515/crelle-2014-0078}.

\bibitem{CGdL} A. Chirre, F. Gon\c{c}alves and D. de Laat,
\emph{Pair correlation estimates for the zeros of the zeta function
via semidefinite programming}, Adv. Math. \textbf{361} (2020),
article 106926,
\href{https://doi.org/10.1016/j.aim.2019.106926}{doi:10.1016/j.aim.2019.106926}.

\bibitem{Devine} M. Devine,
\emph{An Unconditional 67.3399\% Bound and Conditional Advances Beyond
67.92\% for Simple Critical Zeros of the Riemann Zeta Function},
Zenodo preprint, version 1.0.3 (2026),
\href{https://doi.org/10.5281/zenodo.22066689}{doi:10.5281/zenodo.22066689}.

\bibitem{GdLL} F. Gon\c{c}alves, D. de Laat and N. Leijenhorst,
\emph{Multiplicity of nontrivial zeros of primitive $L$-functions
via higher-level correlations}, Math. Comp. \textbf{94} (2025),
no. 354, 2041--2058,
\href{https://doi.org/10.1090/mcom/4005}{doi:10.1090/mcom/4005}.

\bibitem{Johansson} F. Johansson,
\emph{Arb: efficient arbitrary-precision midpoint-radius interval arithmetic},
IEEE Trans. Comput. \textbf{66} (2017), no. 8, 1281--1292,
\href{https://doi.org/10.1109/TC.2017.2690633}{doi:10.1109/TC.2017.2690633}.

\bibitem{Lamzouri} Y. Lamzouri,
\emph{A new proof that more than $2/3$ of the zeros of the Riemann
zeta function are simple and on the critical line},
\href{https://arxiv.org/abs/2609.02882v2}{arXiv:2609.02882v2} (2026).

\bibitem{ZetaLab} T. Lince / Zeta Lab,
\emph{Results: where the seven-point certificate ends}, public research
report, \href{https://github.com/teal-sea/zeta-lab/blob/2da62eb9842db72d4f6bad09c6f13efe384d6ab7/hunts/ainta_seven_point/RESULTS.md}{commit 2da62eb} (2026).

\bibitem{Union} GitHub contributor \texttt{mdumitrean},
public union extension,
\href{https://github.com/anthropics/formal-math/pull/17}{formal-math pull request 17},
accessed 10 September 2026.

\bibitem{Mon73} H. L. Montgomery,
\emph{The pair correlation of zeros of the zeta function},
in Analytic Number Theory (St. Louis, 1972),
Proc. Sympos. Pure Math. \textbf{24}, Amer. Math. Soc., 1973, 181--193,
\href{https://doi.org/10.1090/pspum/024/9944}{doi:10.1090/pspum/024/9944}.

\bibitem{Mon75} H. L. Montgomery,
\emph{Distribution of the zeros of the Riemann zeta function},
in Proceedings of the International Congress of Mathematicians
(Vancouver, 1974), Vol. 1, Canad. Math. Congress, 1975, 379--381.

\bibitem{Ojha} V. Ojha,
\emph{A certified $67.32001\%$ lower-bound candidate for simple zeros
of the Riemann zeta function on the critical line}, public research draft,
\href{https://github.com/trmdy/zeta-simple-zeros-673137/tree/1610b97b7895ff34982260f8dcaf04a0f7b82cf7}{commit 1610b97} (2026).

\bibitem{Shi} Y. Shi,
\emph{A two-certificate trace--energy deduction for simple zeros
of the Riemann zeta function}, public research draft, version 0.1.0,
\href{https://github.com/yuhangshi888/zeta-simple-zeros-673316977/tree/1469eeefe29c48c971ecf98092bc82751dea8bca}{commit 1469eee},
\href{https://doi.org/10.5281/zenodo.21926962}{doi:10.5281/zenodo.21926962} (2026).

\bibitem{STT} A. Simoni\v{c}, T. S. Trudgian and
C. L. Turnage-Butterbaugh,
\emph{Some explicit and unconditional results on gaps between
zeroes of the Riemann zeta-function}, Trans. Amer. Math. Soc.
\textbf{375} (2022), no. 5, 3239--3265,
\href{https://doi.org/10.1090/tran/8571}{doi:10.1090/tran/8571}.

\bibitem{Titchmarsh} E. C. Titchmarsh,
\emph{The theory of the Riemann zeta-function}, second edition,
revised by D. R. Heath-Brown, Clarendon Press, Oxford, 1986.

\bibitem{Waddle} Z. Waddle,
\emph{Unconditional low-multiplicity profiles for zeros of the
Riemann zeta function}, public unrefereed research draft,
\href{https://github.com/zach7036/riemann-hypothesis-research/blob/main/publication/low-multiplicity-zeta/THEOREM.md}{operative theorem statement},
accessed 10 September 2026.

\bibitem{Wang} B. Wang,
\emph{Simple critical zeros and distinct zeros of the Riemann zeta-function
in short intervals},
\href{https://arxiv.org/abs/2609.07918v1}{arXiv:2609.07918v1} (2026).

\end{thebibliography}

\end{document}
